% 论文创新点

本文的主要创新点如下。

（1）设计了基于Markdown结构的知识切分与检索方法。相较于固定长度切分方式，本文按Markdown层级标题进行结构化切分，并结合元数据过滤进行检索，提高了模型在运维场景下获取知识的完整性与准确度。

（2）提出了一种双源运维数据的融合机制。相较于仅依赖单一数据源的方法，本文以故障时间窗为基准，结合TraceID追踪标识，将连续的监控指标与离散的文本日志进行时空维度的对齐关联，构建了结构化的分析上下文，为大模型根因分析提供了良好的数据基础。

（3）实现了一个基于大语言模型的根因分析与指令验证引擎。相较于自由文本生成方式，本文通过结构化Prompt模板规范模型的推理过程与输出格式，并设计了沙盒指令拦截机制，在保证系统安全的前提下，实现了故障修复指令的自动生成与模拟执行验证。
